Looped (weight-tied) Transformers apply a shared residual block $N$ times ($h \leftarrow h + \varepsilon\,f(h)$, same $f$ at each step), increasing effective depth without adding parameters.
Prior depth-scaling analyses prescribe $\varepsilon = 1/\!\sqrt{L}$ for depth-$L$ residual networks.
We show that this is insufficient for looped architectures: weight sharing makes residual updates correlated across iterations, requiring the stronger scaling $\varepsilon = 1/N$.
For multi-layer blocks ($L$ unique layers looped $N$ times), we derive a factored parameterization $\varepsilon = \lambda/(N\!\sqrt{L})$ that separates the two sources of growth: $1/N$ controls the within-layer loop correlation, and $1/\!\sqrt{L}$ controls the across-layer variance.
A key consequence is that the optimal learning rate depends only on the number of unique layers $L$, not on the loop count $N$, enabling direct hyperparameter transfer from small to large $N$ without retuning.
Experiments on looped Transformers confirm that $1/N$ scaling improves trainability and yields better loss than $1/\!\sqrt{N}$ scaling across loop counts.
